%==============================================================================
\chapter{Diseño e implementación del portafolio}
\label{cap:portafolio}
%==============================================================================
% Objetivos específicos 1 y 3. Actividades A1 a A4 y A11 a A15.

Con la base de conocimiento disponible, el problema pasa a ser cómo elegir y
ordenar las configuraciones ante una instancia nueva dentro del presupuesto de
$T = 300$ segundos. Siguiendo la formulación de la
Sección~\ref{sec:formulacion_schedule}, el sistema integra en un mismo flujo la
preparación de la instancia, la extracción de sus características y la
construcción del \textit{schedule} $\sigma(i)$, ya sea mediante reglas de
referencia, métodos de vecinos más cercanos o modelos de aprendizaje
supervisado.

%------------------------------------------------------------------------------
\section{Arquitectura del sistema}
\label{sec:arquitectura}
%------------------------------------------------------------------------------

\subsection{Flujo de una ejecución}

El recorrido de una instancia nueva por el portafolio, desde su recepción hasta
la entrega de la mejor solución, sigue las seis etapas esquematizadas en la
Figura~\ref{fig:flujo_portafolio}:

\begin{enumerate}
  \item En la etapa de preparación, la instancia se convierte al formato
    \texttt{.dzn} (Sección~\ref{sec:instancias}) y se calculan sus cotas $lb$
    y $ub$ (Sección~\ref{sec:cotas}), dejando guardado el programa heurístico
    que produce la $ub$ como solución inicial.
  \item Si en esa verificación inicial se cumple que $lb = ub$, el programa
    heurístico ya es óptimo y el portafolio termina sin ejecutar ningún
    \textit{solver}.
  \item En caso contrario, se calcula el vector de características de la
    instancia y se escala con los mínimos y máximos de la base de
    conocimiento.
  \item A partir de ese vector, una estrategia de selección construye el
    \textit{schedule} definiendo qué configuraciones ejecutar, en qué orden y
    durante cuánto tiempo; en la estrategia principal la decisión se toma con
    tres modelos de aprendizaje automático
    (Sección~\ref{sec:modelos_prediccion}).
  \item A continuación, las configuraciones seleccionadas se ejecutan una tras
    otra con el modelo con cotas, deteniendo el proceso si alguna prueba el
    óptimo o pasando a la siguiente cuando agota su tiempo.
  \item Al concluir la secuencia, se entrega la mejor solución encontrada por
    cualquiera de las configuraciones o, si ninguna mejoró la heurística, el
    programa heurístico.
\end{enumerate}

\begin{figure}[H]
  \centering
  \includegraphics[height=0.55\textheight]{flujo_portafolio.pdf}
  \caption{Flujo de ejecución del portafolio para una instancia nueva.}
  \label{fig:flujo_portafolio}
\end{figure}

\subsection{Componentes}

El sistema se organiza en los componentes de la Tabla~\ref{tab:componentes}.
Cada uno es un programa independiente del repositorio del trabajo
(Anexo~\ref{anx:reproduccion}), de modo que cualquier etapa puede repetirse
por separado. El ejecutor es el componente que construyó la base de
conocimiento (Sección~\ref{sec:ejecucion_solvers}), el cual registra, además
del resultado final, la trayectoria completa de soluciones de cada ejecución
para evaluar cualquier \textit{schedule} con la
ecuación~\eqref{eq:resultado_schedule}.

\begin{table}[H]
  \centering
  \small
  \singlespacing
  \caption{Componentes del sistema y programas que los implementan.}
  \label{tab:componentes}
  \begin{tabularx}{\textwidth}{>{\raggedright\arraybackslash}p{3.4cm} >{\raggedright\arraybackslash}X l}
    \toprule
    \textbf{Componente} & \textbf{Función} & \textbf{Programa (\texttt{src/})} \\
    \midrule
    Conversión de instancias & Traduce el formato de la OR-Library a
      \texttt{.dzn}, con y sin cotas & \texttt{01\_convertir\_instancias.py} \\
    Cálculo de cotas & Algoritmo de Giffler y Thompson con seis reglas y
      verificación de cada programa & \texttt{02\_generar\_cotas.py} \\
    Características genéricas & Extracción con \texttt{mzn2feat} &
      \texttt{03\_caracteristicas\_mzn2feat.py} \\
    Características propias & Características específicas del JSSP &
      \texttt{04\_caracteristicas\_propias.py} \\
    Ejecutor & Lanza cada configuración con MiniZinc, aplica el límite de
      tiempo y registra las trayectorias & \texttt{05\_ejecutor/} \\
    Estrategias de selección & SBS, regla por tamaño, \textit{Winner Takes
      All}, SUNNY y \textit{schedule} estático & \texttt{07\_portafolios.py} \\
    \textit{Schedule} con aprendizaje & Tres modelos de predicción y
      construcción del \textit{schedule} & \texttt{08\_modelo\_ml.py} \\
    \bottomrule
  \end{tabularx}
\end{table}

Para medir el aporte de esta propuesta, el portafolio se compara también con
los mismos \textit{solvers} ejecutados con el modelo original, sin cotas
(Sección~\ref{sec:ejecucion_solvers} y Tabla~\ref{tab:portafolios}).

%------------------------------------------------------------------------------
\section{Composición del portafolio}
\label{sec:composicion}
%------------------------------------------------------------------------------

La primera campaña de ejecución, con el modelo original, incluyó a Chuffed,
Gecode y OR-Tools CP-SAT (Sección~\ref{sec:ejecucion_solvers}). En ella CP-SAT
obtuvo la mejor solución en las 242 instancias, y en 210 de ellas fue el único
en alcanzarla. Dado que esa ventaja se mantiene en casi todo el conjunto, para
estudiar la selección entre alternativas complementarias y con la aprobación de
la dirección del trabajo, el portafolio activo quedó compuesto por tres
configuraciones basadas en Chuffed y Gecode (Tabla~\ref{tab:configuraciones}).

\begin{table}[H]
  \centering
  \small
  \singlespacing
  \caption{Configuraciones del portafolio.}
  \label{tab:configuraciones}
  \begin{tabularx}{\textwidth}{l l >{\raggedright\arraybackslash}X}
    \toprule
    \textbf{Configuración} & \textbf{Modelo} & \textbf{Fortaleza} \\
    \midrule
    Chuffed & Con cotas & Prueba la optimalidad de las instancias pequeñas y
      medianas \\
    Chuffed V2 & Con cotas y estrategia de búsqueda & Encuentra mejores
      soluciones en las instancias grandes \\
    Gecode & Con cotas & Encuentra buenas soluciones en los primeros segundos
      en algunas instancias grandes \\
    \bottomrule
  \end{tabularx}
\end{table}

Chuffed combina búsqueda con aprendizaje de cláusulas y, con su búsqueda
libre, prueba el óptimo de muchas instancias pequeñas y medianas. Chuffed V2 es
el mismo \textit{solver} con la estrategia de búsqueda de la
Sección~\ref{sec:modelo_busqueda}, que prueba menos óptimos pero obtiene
mejores soluciones en las instancias grandes. Gecode, un \textit{solver} de
propagación clásica sin aprendizaje, es el más débil de los tres, pero aporta
diversidad, ya que es un \textit{solver} distinto y en algunas instancias
grandes encuentra la mejor solución en los primeros segundos.
Incluir dos configuraciones de un mismo \textit{solver} es una práctica
habitual en los portafolios basados en configuraciones, como Hydra
\citep{xu2010hydra}, cuyo criterio principal es la complementariedad entre sus
miembros. En la campaña con el modelo con cotas, las tres configuraciones
muestran desempeños complementarios, con ventaja de Chuffed en las instancias
pequeñas y medianas y de Chuffed V2 en las grandes
(Capítulo~\ref{cap:resultados}). CP-SAT se conserva como \textit{solver} de
referencia externo, para situar el desempeño del portafolio respecto a un
\textit{solver} del estado del arte.

%------------------------------------------------------------------------------
\section{Estrategias de selección}
\label{sec:estrategias}
%------------------------------------------------------------------------------

Se implementaron cinco estrategias, de menor a mayor complejidad. Todas usan
únicamente información de la base de conocimiento, que comprende para cada
configuración e instancia de entrenamiento el \textit{makespan} obtenido en
función del tiempo, $C(a, i, \tau)$, y el registro de si probó el óptimo.

\subsection{Mejor \textit{solver} individual}
La estrategia más simple ejecuta siempre la configuración con menor brecha
media en el conjunto de entrenamiento (SBS). En todas las particiones de la
validación resultó ser Chuffed V2. Es la referencia que cualquier otra estrategia
debe superar.

\subsection{Regla por tamaño}
La segunda estrategia usa una sola característica, el número de operaciones
$n \cdot m$. Busca en el entrenamiento el par de configuraciones $(a, b)$ y el
umbral $u$ que minimizan la brecha media de la regla ``ejecutar $b$ si
$n \cdot m \ge u$ y $a$ en otro caso'', y la aplica a la instancia nueva.
Sirve para medir cuánto aporta un modelo con muchas características frente a
una regla que un experto podría escribir a mano.

\subsection{\textit{Winner Takes All} con vecinos más cercanos}
\label{sec:wta}
La estrategia \textit{Winner Takes All} elige una sola configuración y le
asigna los 300 segundos. La elección se hace con el método de los $k$
vecinos más cercanos \citep{cover1967}, escalando las características a
$[-1, 1]$ según los mínimos y máximos del entrenamiento para buscar las $k$
instancias de entrenamiento más cercanas a la nueva según la distancia
euclidiana y elegir la configuración con menor brecha media en ellas. Se usa
$k = 15$, la raíz cuadrada del tamaño del conjunto de entrenamiento
($\sqrt{235} \approx 15$), que es el criterio con el que SUNNY-CP fija el
tamaño del vecindario \citep{amadini2015sunnycp}, y su sensibilidad se
analiza en el Capítulo~\ref{cap:resultados}.

\subsection{\textit{Schedule} secuencial basado en SUNNY}
\label{sec:sunny}
SUNNY \citep{amadini2014sunny, amadini2015sunnycp} construye un
\textit{schedule} para cada instancia a partir de sus vecinos más cercanos.
Su versión original está pensada para problemas de satisfacción, donde una
instancia solo puede quedar resuelta o sin resolver. Como aporte de este
trabajo se implementó una adaptación de SUNNY al JSSP como problema de
optimización, estructurada en cinco pasos:

\begin{enumerate}
  \item \textbf{Vecindario.} Se obtienen los $k = 15$ vecinos más cercanos,
    igual que en \textit{Winner Takes All}.
  \item \textbf{Cobertura.} Se dice que una configuración \textit{cubre} a un
    vecino si alcanza en él el mejor \textit{makespan} del portafolio. Se
    elige el subconjunto de configuraciones que cubre más vecinos y, en caso de
    empate, el de menos configuraciones, seguido del que necesita menos tiempo
    para alcanzar esos valores.
  \item \textbf{Asignación de tiempo.} Cada configuración del subconjunto
    recibe una porción del presupuesto proporcional al número de vecinos que
    cubre.
  \item \textbf{Respaldo.} Los vecinos que ninguna configuración cubre se
    asignan a la configuración con menor brecha media en el vecindario, que
    recibe su porción adicional de tiempo. Este paso proviene del SUNNY
    original; como aquí la cobertura se mide respecto al mejor valor del
    propio portafolio, todo vecino queda cubierto por alguna configuración y
    el paso no llega a activarse, aunque se conserva por completitud.
  \item \textbf{Orden.} Las configuraciones se ejecutan de menor a mayor
    tiempo medio necesario para alcanzar el mejor valor en los vecinos que
    cubren.
\end{enumerate}

\subsection{\textit{Schedule} estático}
\label{sec:schedule_estatico}
Otro aporte realizado en este trabajo es el diseño de un \textit{schedule}
estático con un tramo inicial corto de Gecode, el cual aplica a todas las
instancias el reparto fijo del tiempo con mejor desempeño promedio en el
entrenamiento. Los candidatos evaluados abarcaron todos los pares de
configuraciones en cualquier orden, $[(a, x), (b, T - x)]$, con
$x \in \{2, 5, 10, 15, 20, 30, 45, 60, 90, 120, 150, 200\}$ segundos, los
casos extremos $x = 0$ y $x = 300$ para Gecode y Chuffed, que equivalen a
ejecutar una sola configuración, y tríos que empiezan con 2 segundos de
Gecode, $[(\text{Gecode}, 2), (a, y), (b, T - 2 - y)]$, con
$y \in \{30, 60, 90, 120, 150, 180\}$. Los tríos se incluyeron porque Gecode,
cuando supera a las otras configuraciones, suele hacerlo con las soluciones
que encuentra en los primeros segundos. En todas las particiones se eligió el
mismo \textit{schedule}:
\[
  \sigma_{\mathrm{est}} = \bigl[ (\text{Gecode}, 2), (\text{Chuffed V2}, 60),
  (\text{Chuffed}, 238) \bigr].
\]
Estos valores provienen de una calibración empírica: los tiempos candidatos
se fijaron a mano, cada combinación se evaluó por simulación con las
trayectorias de la base de conocimiento, sin nuevas ejecuciones, y se eligió
la de menor \textit{gap} medio en las instancias de entrenamiento.
Esta estrategia sirve de referencia para medir cuánto aportan las
características de la instancia frente a un reparto fijo bien calibrado.

%------------------------------------------------------------------------------
\section{Modelos de predicción y \textit{schedule} con aprendizaje automático}
\label{sec:modelos_prediccion}
%------------------------------------------------------------------------------

Un aporte central de este trabajo es la estrategia principal del portafolio,
que construye el \textit{schedule} de cada instancia combinando las
predicciones de tres modelos de aprendizaje supervisado. En la literatura de
selección de algoritmos es habitual predecir el desempeño de cada algoritmo a
partir de las características de la instancia \citep{xu2008satzilla,
hutter2014}. Aquí cada modelo responde una de las preguntas que plantea el
\textit{schedule}, determinando si conviene darle tiempo a una configuración
para probar el óptimo, cuánto tiempo necesita y cuál de las configuraciones
dará la mejor solución cuando ninguna lo pruebe.

\subsection{Modelos}
Los tres modelos son bosques aleatorios \citep{breiman2001} de 300 árboles,
con al menos dos ejemplos por hoja, implementados con scikit-learn
\citep{pedregosa2011}. Usan como entrada las características de la instancia
(las de \texttt{mzn2feat} y las propias, escaladas a $[-1, 1]$) y, en los dos
primeros, un indicador de la configuración, de modo que un mismo modelo
aprende el comportamiento de todas.

\begin{itemize}
  \item \textbf{Clasificador de resolución.} Estima la probabilidad
    $p_{\mathrm{res}}(a, i)$ de que la configuración $a$ pruebe el óptimo de
    $i$ dentro de $T$.
  \item \textbf{Regresor de tiempo.} Estima el tiempo $\hat{t}(a, i)$ que
    tardará $a$ en probar el óptimo. Se entrena solo con las ejecuciones que
    lo probaron y sobre el logaritmo del tiempo, porque los tiempos van de
    décimas de segundo a varios minutos. Restringir el entrenamiento a esas
    ejecuciones evita tratar como tiempos reales los valores censurados en el
    límite de 300 segundos, aunque, como advierten \citet{hutter2014},
    descartar los datos censurados tiende a subestimar los tiempos; por eso el
    \textit{schedule} asigna a cada configuración el doble del tiempo predicho.
  \item \textbf{Clasificador de la mejor configuración.} Estima, para cada
    configuración $a$, la probabilidad $p_{\mathrm{mej}}(a, i)$ de que sea la
    que obtiene el menor \textit{makespan} al cabo de los 300 segundos. Si
    varias empatan, se toma como mejor la que alcanza ese valor primero.
\end{itemize}

\subsection{Construcción del \textit{schedule}}
Con las tres predicciones, el \textit{schedule} de una instancia se construye
en dos pasos:

\begin{enumerate}
  \item Primero se ejecutan, de menor a mayor tiempo predicho, las
    configuraciones que probablemente prueban el óptimo, definidas por
    $p_{\mathrm{res}}(a, i) \ge 0.5$. Cada una recibe el doble de su tiempo
    predicho, $2\,\hat{t}(a, i)$, como margen ante el error del regresor, sin
    exceder el tiempo disponible.
  \item A continuación, el tiempo restante $R$ se reparte según la probabilidad
    de que cada configuración sea la mejor. Cada una recibe
    $\max\bigl(m_a,\ p_{\mathrm{mej}}(a, i)\,R\bigr)$ segundos, donde $m_a$ es
    el tiempo que le da el \textit{schedule} estático cuando ese tramo es corto
    (menos de 150 segundos), equivalente a 2 segundos para Gecode y 60 para
    Chuffed V2. Las asignaciones se reescalan para que sumen $R$ y las
    configuraciones se ejecutan de menor a mayor tiempo asignado.
\end{enumerate}

\noindent Así, en una instancia pequeña que Chuffed prueba en pocos segundos,
el \textit{schedule} empieza por Chuffed con un tiempo corto, mientras que en
una instancia grande en la que se espera que gane Chuffed V2, esta
configuración recibe la mayor parte del presupuesto y, cuando el modelo carece
de una predicción clara, el reparto se aproxima al estático.

Todas las estrategias se evalúan con validación \textit{leave-one-out}, de modo
que los modelos, los vecinos, el umbral de la regla por tamaño y el
\textit{schedule} estático se calculan sin la instancia evaluada
(Capítulo~\ref{cap:diseno}).

%------------------------------------------------------------------------------
\section{Entorno de desarrollo y verificación}
\label{sec:entorno}
%------------------------------------------------------------------------------

El sistema se implementó en Python con pandas, NumPy y scikit-learn
\citep{pedregosa2011}, y usa MiniZinc 2.9.1 con Chuffed 0.13.2, Gecode 6.3.0
y OR-Tools CP-SAT 9.12. Para que los resultados fueran confiables, cada
programa construido por la heurística se comprueba contra las restricciones de
precedencia y de capacidad, y se verifica que $lb \le ub$ y que $lb$ no supere
la mejor solución conocida en ninguna instancia. Ningún \textit{makespan}
registrado es menor que la mejor solución conocida de su instancia
\citep{weise2019jssp}, y todo óptimo declarado por un \textit{solver} coincide
con ella. Los programas de cotas y de características propias reproducen
exactamente los archivos usados en el análisis, y todas las tablas y figuras
de resultados se generan con programas del repositorio
(Anexo~\ref{anx:reproduccion}).
